\appendix

\section{Quick Reference Card}
\label{sec:appendix-quickref}

\begin{center}
\small
\begin{tabularx}{\textwidth}{@{}lXXX@{}}
\toprule
\textbf{Property} & \textbf{\pqid{}}                    & \textbf{\pqhash{}}                & \textbf{\pqcrypt{}} \\
\midrule
Variant byte      & \code{0x01} / \code{0x02}           & \code{0x03}                       & \code{0x04} (+ suite byte at offset 7) \\
Wire size         & 41 B fixed                          & 73 B fixed                        & 16-byte header + variable body \\
Tail / body       & 32 B CSPRNG random                  & 64 B SHAKE-256                    & \code{KEM\_CT || AD || nonce || CT || tag || COMMIT} \\
Primitive         & OS CSPRNG                           & SHAKE-256 (FIPS 202)              & ML-KEM-1024 (+ X25519 for \code{0x06}) + AES-256-GCM + UtC \\
Security claim    & 128-bit Grover second-preimage     & $2^{256}$ classical / $\sim 2^{170}$ BHT collision & ML-KEM-1024 IND-CCA2 + AEAD + UtC commitment \\
Public artifact?  & Yes (no secrets)                    & Yes (no secrets)                  & No (carries confidential plaintext) \\
Text encodings    & 55-char Base64URL, 66-char Crockford & 98-char Base64URL, 117-char Crockford & Binary only (v1) \\
Round-trip determinism & N/A (random)                  & Yes (content-deterministic)       & Yes given fixed RNG seed (KAT vectors) \\
\bottomrule
\end{tabularx}
\end{center}

\section{Variant-Byte Allocation Table}
\label{sec:appendix-variants}

\begin{center}
\small
\begin{tabular}{@{}lllll@{}}
\toprule
\textbf{Byte} & \textbf{Owner} & \textbf{Project / variant} & \textbf{First spec version} & \textbf{Status} \\
\midrule
\code{0x00}     & ---        & Permanently reserved (defensive)  & ---     & MUST reject \\
\code{0x01}     & \pqid{}    & Standard generator                & v1.0    & assigned \\
\code{0x02}     & \pqid{}    & HLC generator                     & v1.0    & assigned \\
\code{0x03}     & \pqhash{}  & SHAKE-256 / 64-byte tail          & v1.0    & assigned \\
\code{0x04}     & \pqcrypt{} & Encryption envelope (suite byte at offset 7) & v1.0 & assigned \\
\code{0x05}--\code{0xFE} & --- & Reserved-future                & ---     & MUST reject \\
\code{0xFF}     & ---        & Permanently reserved (defensive)  & ---     & MUST reject \\
\bottomrule
\end{tabular}
\end{center}

\section{PQCrypt Suite-Byte Allocation}
\label{sec:appendix-pqcrypt-suites}

(Header offset 7 within a \pqcrypt{} envelope; only valid when \code{VARIANT == 0x04}.)

\begin{center}
\small
\begin{tabular}{@{}lll@{}}
\toprule
\textbf{Suite} & \textbf{Algorithms}                                                                & \textbf{Status} \\
\midrule
\code{0x01} & ML-KEM-1024 + AES-256-GCM + UtC commit + HKDF-SHA3-256                                & v1.0 mandatory (pure PQ) \\
\code{0x02} & ML-KEM-1024 + ChaCha20-Poly1305 + UtC commit + HKDF-SHA3-256                          & v1.0 optional (pure PQ) \\
\code{0x06} & X25519 + ML-KEM-1024 + AES-256-GCM + UtC commit + HKDF-SHA3-256                       & v1.2 optional (hybrid) \\
\code{0x00}, \code{0x03}--\code{0x05}, \code{0x07}--\code{0xFF} & Reserved-future                & MUST reject \\
\bottomrule
\end{tabular}
\end{center}

\section{Repository Layout}
\label{sec:appendix-layout}

\begin{lstlisting}[language=Pseudocode]
pqtoolkit/
|-- CLAUDE.md                  # umbrella discipline file
|-- README.md                  # repo entry point
|-- latex/                     # this umbrella whitepaper
|   `-- main.tex, sec*.tex, references.bib, Makefile
|
|-- pqid/                      # Project 1: identifier
|   |-- CLAUDE.md, README.md
|   |-- latex/                 # PQID whitepaper sources
|   |-- python/, c/, cpp/, go/, rust/, typescript/, java/, cpython/
|   |-- conformance/vectors/   # KAT files
|   |-- benchmarks/            # per-language bench data
|   `-- docs/                  # release notes, docs site sources
|
|-- pqhash/                    # Project 2: content hash
|   |-- (mirror PQID layout; latex/ added in PQHash v1.0)
|   `-- ...
|
|-- pqcrypt/                   # Project 3: encryption envelope
|   |-- (mirror PQID layout)
|   |-- plugin/                # plugin protocol + reference plugin
|   `-- docs/                  # incl. side-channel-disclaimer.md, security-audit-v1.2.md
|
|-- common/                    # shared library code
|   |-- c/                     # keccak, pqtk_header, random, hkdf_sha3
|   `-- java/                  # Keccak.java
|
`-- tools/
    `-- conformance/           # cross-language KAT harness + per-language runners
\end{lstlisting}

\section{Author and Document Provenance}
\label{sec:appendix-provenance}

\textbf{Author:} Bahram Samani, \pqtoolkit{} monorepo maintainer.

This umbrella whitepaper is built from LaTeX sources under \texttt{latex/} (\texttt{main.tex} + \texttt{secNN\_*.tex}); the build chain is \texttt{pdflatex $\rightarrow$ bibtex $\rightarrow$ pdflatex $\times$ 2} per the Makefile. The artifact at the repo root \texttt{PQToolkit\_Whitepaper.pdf} is regenerated and committed alongside any normative change to the umbrella.

\textbf{Per-project documents.} This umbrella does not restate normative content from the per-project whitepapers. Refer to:

\begin{itemize}
\item \pqid{}: \texttt{pqid/latex/PQID\_Whitepaper.pdf} (repo root: \texttt{PQID\_Whitepaper.pdf}). 46-page document covering wire format, security analysis, distributed-systems model, comparisons with UUIDv4/v7/ULID/KSUID/Snowflake, eight-language reference implementations.
\item \pqhash{}: \texttt{pqhash/latex/PQHash\_Whitepaper.pdf} (repo root: \texttt{PQHash\_Whitepaper.pdf}). 21-page document covering wire format, SHAKE-256 binding, encodings, security argument, conformance, comparison with adjacent hash formats.
\item \pqcrypt{}: \texttt{pqcrypt/latex/PQCrypt\_Whitepaper.pdf} (repo root: \texttt{PQCrypt\_Whitepaper.pdf}). 26-page document covering wire format, ML-KEM-1024 binding, hybrid suite \code{0x06}, AEAD + UtC commitment, KDF, AD-binding patterns, failure modes, side channels, plugin protocol, conformance, comparisons.
\end{itemize}

\textbf{Authority order} (when sources disagree): the per-project whitepaper for project-specific content; this umbrella whitepaper for cross-project content; then the per-project Python reference; then the per-project canonical C ABI; then every other language port. A cross-language disagreement is never resolved by changing the spec to match an outlier.

\textbf{Conformance gate}: \texttt{tools/conformance/run.sh} is the branch-protection check on \texttt{main}. Latest cross-language run (2026-05-21, alpha quick loop): 1051 PASS, 0 FAIL across all three projects on the four alpha languages.

\textbf{Acknowledgments.} Bellare and Hoang for the UtC commitment construction. The OpenSSL, RustCrypto, Go core team, \texttt{kyber-py}, SunJCE, and \texttt{@noble} maintainers for the cryptographic primitive sources \pqcrypt{} vendors. The Cryspen team for \texttt{mlkem-native} (retained as embedded-target fallback). The X-Wing draft authors \citep{xwing-draft} for the combiner design idiom \pqcrypt{}'s hybrid suite borrows. The KyberSlash, PQShield, and GoFetch disclosure teams, whose findings shape \pqcrypt{}'s side-channel disclaimer. The SHAKE-256 / SHA-3 standardization team (Bertoni, Daemen, Peeters, Van Assche) whose primitive is the load-bearing element of every \pqtoolkit{} project.
